\worksheettitle
  {Точки на отрезке}
  {\modulemark{M04} Версия учителя}

\begin{teacherbox}[Паспорт модуля]
  \textbf{Результат:} ученик устанавливает порядок точек, строит схему,
  записывает целое как сумму последовательных частей и находит неизвестную
  длину.

  \textbf{Предварительно:} измеряет отрезки и переводит единицы длины.
\end{teacherbox}

\begin{teacherbox}[Методический акцент]
  Не вводите вычислительную формулу раньше схемы. Сначала ученик проговаривает
  порядок точек и показывает на рисунке целый отрезок и его части.
\end{teacherbox}

\section{Вход и разобранный пример}

Во входе: $CB=11-4=7$ см; целый отрезок $AB$.

\betweenfigure

В примере $CB=13-5=8$ см.

\section{Ответы к практике}

\begin{practicebox}[1. Порядок]
  a) $K$; б) целый $PS$, части $PT$ и $TS$.
\end{practicebox}

\begin{practicebox}[2. Целое и часть]
  \begin{enumerate}[label=\alph*)]
    \item $AB=AC+CB=3\text{ см }8\text{ мм}+5\text{ см }7\text{ мм}
      =9\text{ см }5\text{ мм}$;
    \item $AC=AB-CB=14\text{ см}-6\text{ см }5\text{ мм}
      =7\text{ см }5\text{ мм}$.
  \end{enumerate}
\end{practicebox}

\fourpointsfigure

\begin{practicebox}[Пример с двумя точками]
  $CD=12-4-3=5$ см.
\end{practicebox}

\begin{practicebox}[3. Один общий конец]
  Ближе к $M$ находится $P$.
  \[
    PQ=MQ-MP=8\text{ см }2\text{ мм}-3\text{ см }5\text{ мм}
      =4\text{ см }7\text{ мм}.
  \]
\end{practicebox}

\begin{practicebox}[4. Самостоятельно]
  Порядок $K-M-N-L$,
  \[
    KL=KM+MN+NL,\qquad MN=18-6-5=7\text{ см}.
  \]
\end{practicebox}

\section{Ошибка и контроль}

\begin{warningbox}[Ожидаемое исправление]
  Причина: $AC$ входит в $AD$, потому что оба отрезка начинаются в $A$.
  Поэтому складывать их нельзя. Если порядок $A-C-D$, то
  \[
    CD=AD-AC=9-4=5\text{ см}.
  \]
\end{warningbox}

\begin{checkpointbox}[Ответы]
  Схема должна показывать порядок $A-C-D-B$.
  \[
    AB=AC+CD+DB,\qquad CD=20-7-4=9\text{ см}.
  \]
  Проверка: $7+9+4=20$ см.
\end{checkpointbox}

\begin{teacherbox}[Решение о маршруте]
  M04-R назначается, если ученик не устанавливает порядок, складывает
  вложенные отрезки или не различает целое и части. M04-H достаточно при
  верной модели и единичной вычислительной ошибке. M04\(\star\) не влияет
  на базовый переход.
\end{teacherbox}
